\documentclass[11pt]{article}
\usepackage{mathblatt}
% Präambel des Setzers (werkzeuge/setzer-praeambel.tex, wortgleich); dazu \gteil (Muster T2U)
\makeatletter
\setlength{\mbpfbe}{0mm}% keine Punkte (pruefung.md § 5: nur beim Prüfstein)
\newlength{\gEin}\setlength{\gEin}{0pt}
\let\gSchrift\relax
\renewcommand{\pfkreuzzeile}[1]{\par\noindent\hangindent3.2em\hangafter1\hspace*{1.6em}$\square$\ #1\par}
\newcommand{\gkreuz}[1]{$\square$\ #1\hspace{2em}}
\newcommand{\gkreuzzeile}[2]{\par\noindent\hangindent2.2em\hangafter1\mbox{#1\,$\square$}\ #2\par}% Buchstabe, Option: eine je Zeile (neben Skizze)
% links, was man ansieht, rechts, was man tut; Spaltengrenze überall an derselben Stelle (gemeinsam.md „Satz“)
% v0.17: Fußeintrag „Nr: Wert“ – kurz in einer Box (nie getrennt), lang umbrechbar mit der Nummer am Anfang
% (vorher \mbox: lange Lösungen liefen über den Rand, AS40 138 pt); robust, weil \fusshilfe in die .aux schreibt
\NewDocumentCommand{\gfuss}{m m}{\sbox0{#1:~#2}\ifdim\wd0<0.45\textwidth\mbox{#1:~#2}\else\mbox{#1:}\nobreak\ #2\fi}
\newsavebox{\gbox}\newlength{\gLinks}
% v0.10 (Lehrer 13.10.): keine starre Halbierung – die linke Spalte ist so breit wie die Figur (+ 6 mm), die
% Antwort steht direkt daneben; ist die Figur breiter als 0,62 Zeile, steht die Antwort darunter
\newcommand{\gzwei}[2]{\par\smallskip\noindent\sbox{\gbox}{#1}%
  \setlength{\gLinks}{\dimexpr\wd\gbox+6mm\relax}%
  \ifdim\gLinks>0.62\linewidth
    \usebox{\gbox}\par\smallskip\noindent\begin{minipage}{\linewidth}#2\end{minipage}%
  \else
    \begin{minipage}[t]{\gLinks}\vspace{0pt}\usebox{\gbox}\end{minipage}%
    \begin{minipage}[t]{\dimexpr\linewidth-\gLinks\relax}\vspace{0pt}\raggedright #2\end{minipage}%
  \fi\par}
% v0.9 (12.10., pruefung.md § 5 „Marke mit B, Bedingungen rechts“): drittes Argument von pfaufg = Bedingungen
% der Prüfung (v0.10: \pfohnetr = TR durchgestrichen, „FOR“), klein grau im rechten Rand auf der Grundlinie der Nummer, Zeilen mit \\
\newcommand{\gbedrand}[2]{\ifblank{#2}{}{\rlap{\hspace*{\dimexpr#1+2mm\relax}%
  \parbox[t]{15mm}{\raggedright\scriptsize\leavevmode\color{mbgrau}\strut#2}}}}
% Teil eines Bündels (a, b …): Marke links im Rand, Bedingungen rechts im Rand, auf der Zeile des Teils
\newcommand{\gteilrand}[2]{\llap{\mbpfrandmarke{#1}\hspace*{\mbpfnr}}\gbedrand{\linewidth}{#2}}
% Teil eines Bündels: Buchstabe in eigener Spalte, Inhalt eingerückt (\gEin = Einzug, damit die Spaltengrenze von
% \gzwei an derselben Stelle bleibt)
\newenvironment{gbteil}[3]{\par\smallskip\noindent\gteilrand{#2}{#3}\makebox[\mbpfteil][l]{\textbf{#1}}%
  \begin{minipage}[t]{\dimexpr\linewidth-\mbpfteil\relax}\setlength{\gEin}{\mbpfteil}\raggedright\strut\ignorespaces}%
 {\end{minipage}\par\addvspace{6pt}}% v0.18 (PG80 Nr. 1): Abstand, damit die Linie von a) nicht an b) rutscht
% Grafik höchstens so breit wie die Spalte (Bündel nebeneinander, Koordinatensystem ohne Karo)
\newsavebox{\gpbox}
\newcommand{\gpassend}[1]{\sbox{\gpbox}{#1}\ifdim\wd\gpbox>\linewidth\typeout{GPASSEND: verkleinert \the\wd\gpbox\space auf \the\linewidth}\resizebox{\linewidth}{!}{\usebox{\gpbox}}\else\usebox{\gpbox}\fi}
% Skizzenfeld (v0.10, Lehrer 13.10.): Karo ab dem linken Rand über die Zeile, „Skizze“ klein grau oben links im Feld
% (wie \pfskizzenfeld der Vorlage); der rechte Rand bleibt den Bedingungen
\newcommand{\gskizze}[1]{\par\smallskip\noindent\begin{tikzpicture}
  \clip (0,0) rectangle ({\linewidth-1mm},{#1*10mm});
  \draw[black!16,line width=0.3pt,step=5mm] (0,0) grid ({\linewidth},{#1*10mm});
  \node[anchor=north west,inner sep=2pt,font=\scriptsize,text=mbgrau,fill=white] at (0,{#1*10mm}) {Skizze};
  \end{tikzpicture}\par}
\RenewDocumentEnvironment{pfaufg}{m m m}{%
  \begin{lrbox}{\mbaufgabebox}\begin{minipage}[t]{\linewidth}%
  \gSchrift\noindent\hspace*{\gEin}\mbpfrandmarke{#1}\gbedrand{\linewidth-\gEin}{#3}%
  \begin{minipage}[t]{\mbpfnr}\strut\textbf{#2}\end{minipage}%
  \begin{minipage}[t]{\dimexpr\linewidth-\gEin-\mbpfnr\relax}\raggedright\strut\ignorespaces}%
 {\end{minipage}%
  \end{minipage}\end{lrbox}%
  \setlength{\mbaufgabehoehe}{\dimexpr\ht\mbaufgabebox+\dp\mbaufgabebox\relax}%
  \ifdefempty{\gkopfnach}{\par\addvspace{5pt}\Needspace*{\mbaufgabehoehe}}{%
    % Stufenkopf, der mit dieser (ersten) Aufgabe zusammenbleibt: Platz für beide verlangen, dann den Kopf setzen
    % (12.10., Quadratische: ein Kopf stand allein unten, die Aufgabe mit Koordinatensystem auf der nächsten Seite)
    % v0.17 (AS40 „Terme deuten“ allein unten): hinter dem Kopf steht das Sprungziel (\special, ein Knoten, nach
    % dem die Seite brechen durfte) – darum \nobreak nach dem Kopf; Vorhalt 56 pt = Kopf mit Linie und Abständen
    \par\Needspace*{\dimexpr\mbaufgabehoehe+56pt\relax}\gkopfnach\par\nobreak\gdef\gkopfnach{}\addvspace{5pt}}%
  \noindent\usebox{\mbaufgabebox}\par\addvspace{5pt}}
% v0.18 (plan.md § 10, 15.10. l): Abschnittsbuchstabe A, B, C vor dem Namen im Kopf – als Vorsatz, damit der
% Name im TeX gleich bleibt (Bestellkarte, Sprungziele lesen ihn)
\def\gbstakt{}\newcommand{\gbstvor}[1]{\gdef\gbstakt{#1}}
\NewCommandCopy\gpfstufekopfalt\pfstufekopf
\RenewDocumentCommand{\pfstufekopf}{o m m}{\IfValueTF{#1}{\gpfstufekopfalt[#1]}{\gpfstufekopfalt}%
  {\ifx\gbstakt\empty\else\gbstakt\enspace\fi #2}{#3}\gdef\gbstakt{}}
% v0.18 (16.10. e): Begriff grau, ohne Nummer, ohne Rahmen über der ersten Aufgabe der Stufe
\newcommand{\gbegriff}[1]{\par\nobreak\noindent\hspace*{\mbpfnr}\parbox[t]{\dimexpr\linewidth-\mbpfnr\relax}{\raggedright\small\color{mbgrau}#1}\par\nobreak\addvspace{3pt}}
% v0.18 (16.10. g): Lösungsblatt – Kopf mit kräftigem Strich, Abschnitt klein grau als Zwischenzeile
\newcommand{\glkopf}[1]{\noindent{\large\bfseries Lösungen \textperiodcentered{} #1}\par\vspace{3pt}\hrule height 1.2pt\par\medskip}
\newcommand{\glzwischen}[1]{\par\addvspace{8pt}\Needspace*{4\baselineskip}\noindent{\footnotesize\color{mbgrau}#1}\par\nobreak\vspace{1pt}}
% v0.18 (15.10. o, w): Blattende – Kern ohne Rahmen, dann graue Zeile „Mehr: n Aufgaben“
\newcommand{\gkern}[1]{\par\addvspace{14pt}\Needspace*{6\baselineskip}\noindent{\bfseries Kern}\par\nobreak\smallskip
  {\small\raggedright #1\par}}
\newcommand{\gmehr}[1]{\par\addvspace{10pt}\noindent{\footnotesize\color{mbgrau}Mehr: #1 Aufgaben}\par}
% Stufenkopf vormerken: \gkopfvor{<Kopf>} – die nächste pfaufg setzt ihn, sobald beide zusammen Platz haben
\def\gkopfnach{}
\newcommand{\gkopfvor}[1]{\gdef\gkopfnach{#1}}
% Teilaufgabe mit Skizze links in gzwei: Buchstabe in eigener Spalte, darunter Text und Skizze
\newcommand{\gteil}[2]{\makebox[\mbpfteil][l]{\textbf{#1}}\parbox[t]{\dimexpr0.5\textwidth-\gEin-\mbpfnr-5mm-\mbpfteil\relax}{\raggedright #2}}
% Stufenkopf der Vorlage (mathblatt.sty 2026-10-12: Linie, Abstand, 13 pt, \pfjahre) ohne ihren Vorhalt von
% 14 Zeilen: der Setzer fordert vor dem Kopf 0,16 Seitenhöhe (Kopf + erste Aufgabe), mehr nicht
\setlength{\mbpfstufevorhalt}{0pt}
\makeatother
\newcommand{\karo}[1]{\par\smallskip\noindent\begin{tikzpicture}
  \clip (0,0) rectangle ({\linewidth-1mm},{#1*10mm});
  \draw[black!16,line width=0.3pt,step=5mm] (0,0) grid ({\linewidth},{#1*10mm});
  \end{tikzpicture}\par}
\newcommand{\kk}[1]{$\square$\,#1\hspace{0.9em}}
\newcommand{\linie}[1]{\,{\color{black!45}\rule[-1pt]{#1}{0.4pt}}\,}
\newcommand{\pkt}{\textperiodcentered{}}
% v0.14 (14.10.): Stufenkopf mit fünf Punkten (Option = Jahre für die Bestellkarte) oder „zuletzt 19“
\NewDocumentCommand{\gpunkte}{o m}{#2}
\newcommand{\gzuletzt}[1]{zuletzt #1}
% v0.7 (12.10.): Schülerfrage unter dem Stufenkopf; Ankreuzoption als Bild (Buchstabe und Kästchen darunter)
\newcommand{\gfrage}[1]{\par\nobreak\noindent\hspace*{\mbpfnr}{\small\itshape #1}\par\nobreak\addvspace{2pt}}
\newlength{\gbildbreite}\setlength{\gbildbreite}{0.32\linewidth}% Setzer setzt sie je Reihe (vier Bilder: 0,24)
\newcommand{\gbildoption}[3]{\begin{minipage}[t][#3][b]{\gbildbreite}\centering #2\par\smallskip\mbox{#1\,$\square$}\end{minipage}}
% Lineare (12.10.): Bild mit Buchstabe ohne Kästchen (Diagramm zuordnen); Aussage mit „wahr / falsch“ zum Ankreuzen
\newcommand{\gbildname}[3]{\begin{minipage}[t][#3][b]{\gbildbreite}\centering #2\par\smallskip\mbox{\textbf{#1}}\end{minipage}}
\newcommand{\gwahrfalsch}[2]{\par\noindent\hangindent2.2em\hangafter1\makebox[2.2em][l]{(#1)}#2\hfill\mbox{wahr $\square$\quad falsch $\square$}\par}
\begin{document}

\pfheftstil
\pflaufkopf[Abitur GK \pkt{} Prüfung]{Graph im Sachzusammenhang \pkt{} Lösungen}{AA130}\pfseitevon
\pfheftkopf{Graph im Sachzusammenhang}{Lösungen}
\pfstufenzeile{Graph lesen und skizzieren \pkt{} Länge gesucht \pkt{} Maß mit Maßstab gesucht}

\Needspace*{0.16\textheight}\pfstufekopf{Graph lesen und skizzieren}{\gpunkte[25 23]{\textcolor{black!25}{$\circ$}\,\textcolor{black!80}{$\bullet$}\,\textcolor{black!25}{$\circ$}\,\textcolor{black!80}{$\bullet$}\,\textcolor{black!25}{$\circ$}}}\special{pdf:dest (stufe-graph-lesen-und-skizzieren) [@thispage /XYZ @xpos @ypos null]}
\begin{pfloesungb}{}
\lzb{1.}{}{in der Dauerphase schwankt die Temperatur zwischen etwa 181 °C und 208 °C; Durchschnitt $\approx$ 195 °C (190 °C bis 200 °C vertretbar)}{in der Dauerphase schwankt die Temperatur zwischen etwa 181 °C und 208 °C; Durchschnitt $\approx$ 195 °C (190 °C bis 200 °C vertretbar)}
\end{pfloesungb}
\begin{pfloesungb}{}
\lzb{2.}{}{die Zahl der Likes nimmt zunächst immer stärker zu (bis etwa 13:25 Uhr, Wendestelle $\approx$ 1,4); danach nimmt sie immer schwächer zu und nähert sich 2400}{die Zahl der Likes nimmt zunächst immer stärker zu (bis etwa 13:25 Uhr, Wendestelle $\approx$ 1,4); danach nimmt sie immer schwächer zu und nähert sich 2400}
\end{pfloesungb}
\begin{pfloesungb}{}
\lzb{3.}{}{Skizze durch (0|$-$16), H(2|4), W(3|2), T(4|0), (6|20); falsch}{Skizze durch (0|$-$16), H(2|4), W(3|2), T(4|0), (6|20); falsch: die Wendetangente y = $-$3x + 11 berührt den Graphen nur in W, links und rechts liegt der Graph auf verschiedenen Seiten der Tangente}
\end{pfloesungb}

\Needspace*{0.16\textheight}\pfstufekopf{Länge gesucht}{\gpunkte[26 22]{\textcolor{black!80}{$\bullet$}\,\textcolor{black!25}{$\circ$}\,\textcolor{black!25}{$\circ$}\,\textcolor{black!25}{$\circ$}\,\textcolor{black!80}{$\bullet$}}}\special{pdf:dest (stufe-lange-gesucht) [@thispage /XYZ @xpos @ypos null]}
\begin{pfloesungb}{}
\lzb{4.}{}{Höhe 2 $-$ ($-$2) = 4; A = 4$\pi$}{Grundseite AC = 5$\pi$/2 $-$ $\pi$/2 = 2$\pi$; Höhe 2 $-$ ($-$2) = 4; A = $\tfrac{1}{2}$ · 2$\pi$ · 4 = 4$\pi$}
\end{pfloesungb}
\begin{pfloesungb}{}
\lzb{5.}{}{Abstieg bis t = 12: 864 m; 864 $-$ 729 = 135 m; Weg 864 + 135 = 999 m}{Abstieg bis t = 12: 864 m; h(18) = 729 $\Rightarrow$ Aufstieg von 12 bis 18 min: 864 $-$ 729 = 135 m; Weg 864 + 135 = 999 m}
\end{pfloesungb}
\begin{pfloesungb}{}
\lzb{6.}{a)}{Halbkreisradius $\tfrac{4}{3}$; Höhe 2 · $\tfrac{4}{3}$ $\approx$ 2,67 m $\le$ 2,7 m; passt}{Halbkreisradius $\tfrac{4}{3}$ (Abstand von A zur x-Achse); Höhe 2 · $\tfrac{4}{3}$ = $\tfrac{8}{3}$ m $\approx$ 2,67 m $\le$ 2,7 m; Breite 2 · (1 + $\tfrac{4}{3}$) = $\tfrac{14}{3}$ m $\approx$ 4,67 m $\le$ 4,7 m $\Rightarrow$ passt}
\end{pfloesungb}
\begin{pfloesungb}{}
\lzb{}{b)}{4 · 1,7 = 6,8 m; 2$\pi$ · $\tfrac{4}{3}$ $\approx$ 8,38 m; Zeit $\tfrac{15{,}18}{3}$ $\approx$ 5 min 4 s}{vier Bögen: 4 · 1,7 = 6,8 m; zwei Halbkreise = ein Vollkreis: 2$\pi$ · $\tfrac{4}{3}$ $\approx$ 8,38 m; Länge $\approx$ 15,18 m $\Rightarrow$ Zeit $\tfrac{15{,}18}{3}$ $\approx$ 5,06 min $\approx$ 5 min 4 s}
\end{pfloesungb}

\Needspace*{0.16\textheight}\pfstufekopf{Maß mit Maßstab gesucht}{\gpunkte[24 23]{\textcolor{black!25}{$\circ$}\,\textcolor{black!25}{$\circ$}\,\textcolor{black!80}{$\bullet$}\,\textcolor{black!80}{$\bullet$}\,\textcolor{black!25}{$\circ$}}}\special{pdf:dest (stufe-mass-mit-massstab-gesucht) [@thispage /XYZ @xpos @ypos null]}
\begin{pfloesungb}{}
\lzb{7.}{}{quadratische Grundfläche mindestens 60 cm × 60 cm; Höhe 60 cm; V = 216 l}{Platte 6 LE × 6 LE = 60 cm × 60 cm $\Rightarrow$ quadratische Grundfläche mindestens 60 cm × 60 cm; Höhe 60 cm; V = 60 · 60 · 60 cm³ = 216 000 cm³ = 216 l}
\end{pfloesungb}
\begin{pfloesungb}{}
\lzb{8.}{a)}{Radius f(12) $\approx$ 1,218 LE; $\approx$ 2,436 LE; · 0,5 cm $\approx$ 12,2 mm}{Radius f(12) = 9 · $e^{-2}$ $\approx$ 1,218 LE; Durchmesser 2 · 1,218 LE $\approx$ 2,436 LE; · 0,5 cm $\approx$ 1,22 cm $\approx$ 12,2 mm}
\end{pfloesungb}
\begin{pfloesungb}{}
\lzb{}{b)}{Länge = 4,5 cm; Breite = 3 cm; Masse 40,5 · 0,86 $\approx$ 34,8 g}{12 $-$ 3 = 9 LE $\Rightarrow$ Länge = 4,5 cm; 2 · 3 = 6 LE $\Rightarrow$ Breite = Höhe = 3 cm; V = 4,5 · 3 · 3 = 40,5 cm³ $\Rightarrow$ Masse 40,5 · 0,86 $\approx$ 34,8 g}
\end{pfloesungb}
\end{document}